\section{The source fragment \Frag}
\label{sec:fragment}

We now define the source calculus: a polymorphic first-order fragment of CIC
with parameterized inductive data types and inductive predicates.  The
fragment is chosen to contain the phenomena that make guard and
type-argument omission non-trivial --- polymorphism, dependent inductive
predicates (predicates whose well-formedness constrains their arguments'
types), empty types, equality --- while excluding the CIC features
(higher-order quantification, dependently typed \emph{data}, conversion)
whose treatment by hammer translations involves further, orthogonal
machinery; \cref{rem:scope} delimits the boundary precisely.

\subsection{Syntax and typing}

\begin{definition}[Types]
  \label{def:types}
  Fix a countably infinite set of \emph{type variables}
  $\alpha, \beta, \dots$.  Given a set $\Ind$ of \emph{inductive type
  formers}, each $I \in \Ind$ with an arity $k_I \geq 0$, the \emph{types
  over a set $V$ of type variables} are generated by
  \[
    \tau ::= \alpha \mid I(\seq{\tau}{k_I}), \qquad \alpha \in V,\ I \in \Ind.
  \]
  We write $\Ty{V}$ for this set, $\ftv{\tau}$ for the type variables
  occurring in $\tau$, and $\GTy \defeq \Ty{\emptyset}$ for the set of
  \emph{ground types}.  A \emph{type substitution} is a finite map
  $\theta$ from type variables to types, applied in the usual way and
  written postfix: $\tau\theta$.
\end{definition}

\begin{definition}[Environments]
  \label{def:environment}
  An \emph{environment} $\Env$ consists of:
  \begin{enumerate}[label=(\alph*)]
  \item a finite set $\Ind$ of inductive type formers with arities;
  \item a finite set $\Constr$ of \emph{constructors}; each $c \in \Constr$
    is assigned to exactly one $I \in \Ind$ and declared as
    \[
      c : \forall \seq{\alpha}{k_I}.\
          \tau^c_1 \times \dots \times \tau^c_{m_c} \arr I(\seq{\alpha}{k_I}),
      \qquad \tau^c_i \in \Ty{\{\seq{\alpha}{k_I}\}};
    \]
  \item a finite set $\Fun$ of \emph{function constants}, each declared as
    $f : \forall \seq{\alpha}{k_f}.\ \tau^f_1 \times \dots \times
    \tau^f_{m_f} \arr \tau^f_0$ with all $\tau^f_i \in
    \Ty{\{\seq{\alpha}{k_f}\}}$;
  \item a finite set $\Pred$ of \emph{predicate constants}, each declared as
    $p : \forall \seq{\alpha}{k_p}.\ \tau^p_1 \times \dots \times
    \tau^p_{m_p} \arr \PProp$ with all $\tau^p_i \in
    \Ty{\{\seq{\alpha}{k_p}\}}$;
  \item a finite set $\mathrm{Prem}(\Env)$ of closed formulas
    (\cref{def:formulas}), the \emph{premises}.
  \end{enumerate}
  Zero-ary cases are allowed throughout: $I$ may have no constructors (an
  empty type), $c$ or $f$ may have no term arguments, and any of the symbol
  classes may be empty.
\end{definition}

\begin{definition}[Terms and typing]
  \label{def:terms}
  A \emph{context} is a pair $(V;\Gamma)$ of a finite set $V$ of type
  variables and a finite map $\Gamma$ from term variables to $\Ty{V}$.
  Terms are generated by
  \[
    t ::= x \mid h\langle\seq{\sigma}{k_h}\rangle(\seq{t}{m_h}),
    \qquad h \in \Constr \cup \Fun,
  \]
  where constants are always fully applied to $k_h$ type arguments and
  $m_h$ term arguments.  The typing judgment $\wfj{V}{\Gamma}{t : \tau}$
  (with $\tau \in \Ty{V}$) is defined by two syntax-directed rules:
  \begin{align*}
    &\frac{\Gamma(x) = \tau}{\wfj{V}{\Gamma}{x : \tau}}
    &&\frac{
      \sigma_1,\dots,\sigma_{k_h} \in \Ty{V} \qquad
      \wfj{V}{\Gamma}{t_i : \tau^h_i[\tup\sigma/\tup\alpha]}
      \quad (1 \le i \le m_h)
    }{
      \wfj{V}{\Gamma}{h\langle\tup\sigma\rangle(\tup t) :
        \tau^h_0[\tup\sigma/\tup\alpha]}
    }
  \end{align*}
  where for $h = c \in \Constr$ assigned to $I$ we set
  $\tau^c_0 \defeq I(\seq{\alpha}{k_I})$.
\end{definition}

\begin{lemma}[Uniqueness of types]
  \label{lem:unique-types}
  If $\wfj{V}{\Gamma}{t : \tau}$ and $\wfj{V}{\Gamma}{t : \tau'}$ then
  $\tau = \tau'$.
\end{lemma}

\begin{proof}
  Immediate induction on $t$: the typing rules are syntax-directed and the
  type in the conclusion is a function of the syntax ($\Gamma(x)$ for a
  variable; $\tau^h_0[\tup\sigma/\tup\alpha]$ for an application, where
  $\tup\sigma$ is written in the term).
\end{proof}

\begin{lemma}[Stability under instantiation]
  \label{lem:instantiation}
  Let $\theta$ be a type substitution with $\alpha\theta \in \Ty{V'}$ for
  all $\alpha \in V$.  If $\wfj{V}{\Gamma}{t : \tau}$ then
  $\wfj{V'}{\Gamma\theta}{t\theta : \tau\theta}$, where $\theta$ acts on
  terms by replacing each type-argument tuple
  $\langle\tup\sigma\rangle$ by $\langle\tup\sigma\theta\rangle$.
\end{lemma}

\begin{proof}
  Induction on the typing derivation, using
  $\tau^h_i[\tup\sigma/\tup\alpha]\theta =
  \tau^h_i[\tup\sigma\theta/\tup\alpha]$, which holds because
  $\ftv{\tau^h_i} \subseteq \{\tup\alpha\}$ and substitution composes.
\end{proof}

\begin{definition}[Formulas]
  \label{def:formulas}
  Formulas and their well-formedness judgment $\wfj{V}{\Gamma}{\phi\
  \mathrm{wf}}$ are generated by:
  \[
    \phi ::= p\langle\tup\sigma\rangle(\tup t)
    \mid t \eqx{\tau} s \mid \bot \mid \top
    \mid \phi \wedge \phi \mid \phi \vee \phi \mid \phi \imp \phi
    \mid \forall x{:}\tau.\ \phi \mid \exists x{:}\tau.\ \phi
    \mid \forall \alpha.\ \phi
  \]
  with the evident rules: an atom $p\langle\tup\sigma\rangle(\tup t)$ is
  well formed when $\tup\sigma \in \Ty{V}^{k_p}$ and
  $\wfj{V}{\Gamma}{t_i : \tau^p_i[\tup\sigma/\tup\alpha]}$ for all $i$; an
  equation $t \eqx{\tau} s$ when both sides have type $\tau \in \Ty{V}$;
  quantifiers extend $\Gamma$ (resp.\ $V$).  We use
  $\neg\phi \defeq \phi \imp \bot$ and $\biimp$ as abbreviations.  A
  \emph{sentence} is a formula well formed in the empty context; premises
  in $\mathrm{Prem}(\Env)$ are required to be sentences.  \emph{Type
  quantification is prenex}: in $\forall\alpha.\phi$ we require that
  $\phi$ contain no term quantifier scoping over a later type quantifier;
  equivalently, every sentence has the form
  $\forall \seq{\alpha}{n}.\ \psi$ with $\psi$ free of type quantifiers.
  \Cref{lem:instantiation} extends to formulas by the same induction.
\end{definition}

\subsection{The ground companion theory}
\label{subsec:companion}

The translation's soundness will be stated relative to a many-sorted
first-order theory that collects everything CIC proves about
\Frag-sentences by case analysis on inductive data and inductive proofs.
We first fix the many-sorted setting.

\begin{definition}[Ground signature and structures]
  \label{def:ms-structures}
  The \emph{ground signature} of $\Env$ has sort set $\GTy$ and, for every
  $h \in \Constr \cup \Fun$ and every $\tup\tau \in \GTy^{k_h}$, a function
  symbol $h_{\tup\tau}$ of rank
  $\tau^h_1[\tup\tau/\tup\alpha] \times \dots \times
  \tau^h_{m_h}[\tup\tau/\tup\alpha] \arr \tau^h_0[\tup\tau/\tup\alpha]$,
  and for every $p \in \Pred$, $\tup\tau \in \GTy^{k_p}$ a predicate symbol
  $p_{\tup\tau}$ of rank
  $\tau^p_1[\tup\tau/\tup\alpha] \times \dots \times
  \tau^p_{m_p}[\tup\tau/\tup\alpha]$.  A \emph{ground structure} $\NN$
  assigns to every sort $\tau \in \GTy$ a (possibly empty) carrier set
  $\NN_\tau$, to every $h_{\tup\tau}$ a function between the corresponding
  carriers, and to every $p_{\tup\tau}$ a relation; each equality
  $\eqx{\tau}$ is interpreted as identity on $\NN_\tau$.  Satisfaction
  $\NN, v \models \phi$ for type-quantifier-free formulas $\phi$ with a
  sort-respecting valuation $v$ is defined as usual (a universal
  (existential) quantifier over an empty carrier is true (false)).  For a
  sentence with prenex type quantifiers we define \emph{ground
  satisfaction}
  \[
    \NN \entg \forall\seq{\alpha}{n}.\ \psi
    \quad\text{iff}\quad
    \NN \models \psi[\tup\tau/\tup\alpha]
    \ \text{for every } \tup\tau \in \GTy^n,
  \]
  where the ground instance $\psi[\tup\tau/\tup\alpha]$ is read over the
  ground signature by interpreting each
  $h\langle\tup\sigma\rangle$ as $h_{\tup\sigma[\tup\tau/\tup\alpha]}$
  (well-typedness of instances is \cref{lem:instantiation}).  We allow
  empty carriers deliberately: an inductive type former with no
  constructors, or with constructors whose argument types are empty, denotes
  an empty collection, and the exhaustiveness axioms below force this.
\end{definition}

\begin{definition}[Structural axioms]
  \label{def:struct}
  $\StructAx$ is the following set of \Frag-sentences.
  For every $I \in \Ind$ with constructors $c^{(1)},\dots,c^{(r)}$ (possibly
  $r=0$) and for $\tup\alpha$ of length $k_I$:
  \begin{enumerate}[label=(S\arabic*)]
  \item \emph{Injectivity}, for each constructor $c = c^{(j)}$:
    \[
      \forall\tup\alpha.\ \forall \tup x{:}\tup{\tau^c}.\ \forall
      \tup x'{:}\tup{\tau^c}.\
      c\langle\tup\alpha\rangle(\tup x) \eqx{I(\tup\alpha)}
      c\langle\tup\alpha\rangle(\tup x')
      \imp \textstyle\bigwedge_i x_i \eqx{\tau^c_i} x'_i .
    \]
  \item \emph{Distinctness}, for each pair $j \neq j'$:
    $\forall\tup\alpha\,\tup x\,\tup x'.\ \neg\bigl(
    c^{(j)}\langle\tup\alpha\rangle(\tup x) \eqx{I(\tup\alpha)}
    c^{(j')}\langle\tup\alpha\rangle(\tup x')\bigr)$.
  \item \emph{Exhaustiveness}:
    $\forall\tup\alpha.\ \forall z{:}I(\tup\alpha).\
    \bigvee_{j=1}^{r} \exists \tup x{:}\tup{\tau^{c^{(j)}}}.\
    z \eqx{I(\tup\alpha)} c^{(j)}\langle\tup\alpha\rangle(\tup x)$
    \quad (the empty disjunction being $\bot$).
  \end{enumerate}
  If some predicates of $\Env$ are designated as \emph{inductively defined}
  with given clauses, $\StructAx$ additionally contains their introduction
  rules and their inversion principles, each of which is a \Frag-sentence
  (for $\mathit{le}$, the inversion principle is the formula displayed in
  \cref{sec:overview} before translation).  We do not fix a clause format;
  we only require that whatever is put into $\StructAx$ for predicates be
  CIC-provable in the sense of \cref{prop:cic-bridge}.
\end{definition}

\begin{definition}[Ground companion theory]
  \label{def:companion}
  $\Th$ is the theory $\mathrm{Prem}(\Env) \cup \StructAx$, considered via
  ground satisfaction: $\NN \models \Th$ means $\NN \entg \phi$ for every
  $\phi \in \mathrm{Prem}(\Env) \cup \StructAx$.  We write
  $\Th \entg \phi_0$ to mean that every ground structure satisfying $\Th$
  ground-satisfies the sentence $\phi_0$.  $\Env$ is \emph{consistent} if
  $\Th$ has at least one ground model.
\end{definition}

\subsection{Relation to CIC}
\label{subsec:cic}

The fragment embeds into CIC in the obvious way: each $I$ becomes a
parameterized inductive type in $\mathsf{Type}$, constructors become
constructors, function constants become section variables (or definitions)
of the corresponding simple type, predicate constants become variables of
predicate type or inductive definitions in $\PProp$, formulas become
propositions --- with $\eqx{\tau}$ read as Leibniz equality and
$\forall\alpha.\phi$ as $\Pi \alpha{:}\mathsf{Type}.\,\phi$ --- and
premises become hypotheses.  Call $\Env$ \emph{CIC-realizable} if the
embedded environment is well formed and every premise's embedding is
inhabited.

\begin{proposition}[CIC bridge]
  \label{prop:cic-bridge}
  Assume $\Env$ is CIC-realizable and the embeddings of the structural
  axioms (S1)--(S3) and of the predicate introduction/inversion sentences
  are inhabited --- which they are whenever the embedded predicates are
  inductively defined with the corresponding clauses, since (S1) and (S2)
  are proved by dependent case analysis and an injection/discrimination
  argument, (S3) by case analysis on $z$, and predicate inversion by case
  analysis on the proof.  If CIC extended with excluded middle is
  consistent, then $\Env$ is consistent in the sense of
  \cref{def:companion}.
\end{proposition}

\begin{proof}
  Interpret the embedded environment in the set-theoretic model of CIC with
  classical logic \cite{Werner1997,LeeWerner2011} (any model of CIC+EM
  suffices).  Define $\NN$ by letting $\NN_\tau$ be the set interpreting the
  embedding of the ground type $\tau$, and interpreting each $h_{\tup\tau}$,
  $p_{\tup\tau}$ by the interpretation of the embedded constant at the
  corresponding type instance.  A straightforward induction on
  type-quantifier-free \Frag-formulas shows that classical set-theoretic
  truth of the embedding of a ground instance coincides with
  $\NN \models \cdot$ of that instance: atoms and equations agree by
  construction, the connectives are interpreted classically on both sides,
  and term quantifiers range over $\NN_\tau$ on both sides.  Every sentence
  of $\Th$ is inhabited in CIC(+EM) by assumption, hence true in the model,
  hence all its ground instances are true, hence $\NN \entg$ each of them.
  Thus $\NN \models \Th$.
\end{proof}

From here on we never mention CIC again until \cref{sec:related}: all
theorems are stated for an arbitrary environment $\Env$ and quantify over
ground models $\NN \models \Th$.  \Cref{prop:cic-bridge} is the (only)
bridge: for a CIC-realizable environment, ``$\Th \entg \phi_0$'' is a sound
notion of the goal $\phi_0$ being a (classical) consequence of the
development, and a conclusion of the form ``the optimized translation is
refutable $\Rightarrow \Th \entg \phi_0$'' is exactly the guarantee that
the ATP has not proved anything beyond what the source environment
classically entails.

\begin{remark}[Scope]
  \label{rem:scope}
  Compared to full CIC, \Frag\ excludes: function types as data (arguments
  and results of constants are inductive types or type variables, so the
  translation needs no lambda-lifting and no extensionality reasoning);
  dependently typed data (type formers take only type parameters, not term
  indices --- but \emph{predicates} may have term arguments of arbitrary
  fragment types, which is where the interesting covers arise);
  higher-rank polymorphism (type quantification is prenex, as in ML);
  universes above the single implicit $\mathsf{Type}$; and conversion
  (constants are opaque; definitional equations enter as premises).  These
  restrictions match the first-order core that hammer translations handle
  natively; the excluded features are handled by upstream devices
  (lifting, collapsing) whose interaction with omission we discuss in
  \cref{sec:related}.  Within the fragment, however, we impose no
  restriction on the shapes of premises, on emptiness of types, or on which
  guards and type arguments occur where.
\end{remark}
